\chapter{Overall Conclusion}

The GIPEDI Internship at the Department of Electrical Engineering, Indian Institute of Technology Delhi provided valuable hands-on experience in embedded systems, electronic hardware development, Arduino programming, and PCB design. Throughout the internship, theoretical concepts learned during academic coursework were applied to practical engineering problems, strengthening both technical knowledge and problem-solving skills.

During the internship, several embedded system experiments and applications were studied and implemented using the Arduino platform. These included sensor interfacing, display modules, real-time clock applications, EEPROM data logging, motor control, waveform generation, weather monitoring, security systems, and other microcontroller-based experiments. These activities strengthened my understanding of embedded programming, digital and analog interfacing, communication protocols, and hardware debugging while providing a strong practical foundation in embedded system design.

The first major project focused on the development of a load cell-based smart weighing system using a 1~kg load cell, an HX711 precision amplifier, an ESP32 Dev Module, and a SH1106 OLED display. The system was successfully designed, assembled, calibrated, and experimentally evaluated. The sensitivity analysis demonstrated that the developed prototype could reliably detect weight variations of approximately 1~g while providing stable and repeatable measurements. This project provided practical experience in sensor interfacing, calibration techniques, embedded programming, hardware integration, and experimental data analysis.

The second major project involved the design of a low-cost White Noise Generator Printed Circuit Board (PCB) using KiCad. The complete workflow, including schematic capture, footprint assignment, component placement, copper trace routing, design verification, and generation of manufacturing outputs, was successfully completed by following standard PCB design practices. Although PCB fabrication and hardware testing were outside the scope of the internship, the final PCB design serves as a manufacturable prototype suitable for future implementation.

Overall, this internship significantly enhanced my practical skills in embedded system development, electronic circuit design, PCB layout design, debugging, technical documentation, and engineering problem solving. It also provided valuable exposure to professional engineering tools and technologies such as Arduino IDE, ESP32, KiCad, SimulIDE, sensors, communication modules, and modern PCB design methodologies. The knowledge and experience gained during this internship have established a strong foundation for future academic research, industrial projects, and the development of advanced embedded and Internet of Things (IoT) systems.